\chapter{The Standard Discrete Distributions}
\label{ch:dist}

A small family of distributions accounts for most discrete models. For each we
state the model, derive the PMF, and compute the mean and variance. We use only
tools developed so far; where independence of random variables would give a
shorter derivation, we note it and give the proof in Chapter~\ref{ch:joint}.

\section{Bernoulli and binomial}

\begin{definition}[Bernoulli]
$X$ has the \emph{Bernoulli}$(p)$ distribution if $\Prob(X=1)=p$ and
$\Prob(X=0)=1-p=q$. This is the indicator of a ``success'' in one trial.
\end{definition}

Its mean is $p$ and its variance is $pq$ (Chapter~\ref{ch:drv}, indicators).

\begin{definition}[Binomial]
$X$ has the \emph{binomial}$(n,p)$ distribution if
\[
\Prob(X=k)=\binom nkp^kq^{n-k},\qquad k=0,1,\dots,n .
\]
It is the number of successes in $n$ independent trials, each with success
probability $p$ (Example~\ref{ex:trials}).
\end{definition}

The probabilities sum to $1$ by the binomial theorem:
$\sum_k\binom nkp^kq^{n-k}=(p+q)^n=1$.

\begin{theorem}[Mean and variance of the binomial]\label{thm:binmoments}
If $X\sim\mathrm{Bin}(n,p)$ then $\Exp[X]=np$ and $\Var(X)=npq$.
\end{theorem}
\begin{proof}
We use the identities $k\binom nk=n\binom{n-1}{k-1}$ and
$k(k-1)\binom nk=n(n-1)\binom{n-2}{k-2}$, both immediate from the factorial
formula. Then
\[
\Exp[X]=\sum_{k=1}^nk\binom nkp^kq^{n-k}=np\sum_{k=1}^n\binom{n-1}{k-1}p^{k-1}q^{n-k}
=np\,(p+q)^{n-1}=np.
\]
Similarly
\[
\Exp[X(X-1)]=\sum_{k=2}^nk(k-1)\binom nkp^kq^{n-k}=n(n-1)p^2(p+q)^{n-2}=n(n-1)p^2.
\]
Hence $\Exp[X^2]=\Exp[X(X-1)]+\Exp[X]=n(n-1)p^2+np$ and
\[
\Var(X)=n(n-1)p^2+np-n^2p^2=np-np^2=npq.\qedhere
\]
\end{proof}

(Probabilistic shortcut: $X$ is a sum of $n$ independent Bernoulli$(p)$ variables;
linearity gives the mean at once and independence gives the variance, as shown in
Chapter~\ref{ch:joint}.)

\begin{proposition}[Mode of the binomial]
The ratio $\dfrac{\Prob(X=k)}{\Prob(X=k-1)}=\dfrac{(n-k+1)p}{kq}$ exceeds $1$ iff
$k<(n+1)p$. So the PMF increases up to $\lfloor(n+1)p\rfloor$ and decreases
afterwards; the mode is $\lfloor(n+1)p\rfloor$ (with a tie at $k=(n+1)p$ when this
is an integer).
\end{proposition}
\begin{proof}
The ratio is a direct computation from the PMF. The inequality
$(n-k+1)p>kq$ rearranges to $(n+1)p>k(p+q)=k$.
\end{proof}

\section{Geometric and negative binomial}

\begin{definition}[Geometric]
$X$ has the \emph{geometric}$(p)$ distribution, $0<p\le1$, if
\[
\Prob(X=k)=q^{k-1}p,\qquad k=1,2,\dots
\]
It is the number of independent trials up to and including the first success.
\end{definition}

The total mass is $\sum_{k\ge1}q^{k-1}p=p/(1-q)=1$. The tail is
$\Prob(X\ge k)=q^{k-1}$ (the first $k-1$ trials fail).

\begin{theorem}[Moments of the geometric]\label{thm:geomoments}
$\Exp[X]=1/p$ and $\Var(X)=q/p^2$.
\end{theorem}
\begin{proof}
By the tail-sum formula (Proposition~\ref{prop:expprops}(v)),
$\Exp[X]=\sum_{k\ge1}q^{k-1}=1/(1-q)=1/p$.
For the variance compute $\Exp[X(X-1)]=\sum_{k\ge2}k(k-1)q^{k-1}p$. Differentiating
the geometric series $\sum_{k\ge0}x^k=(1-x)^{-1}$ twice (valid for $\abs x<1$)
gives $\sum_{k\ge2}k(k-1)x^{k-2}=2(1-x)^{-3}$. Hence
\[
\Exp[X(X-1)]=pq\sum_{k\ge2}k(k-1)q^{k-2}=\frac{2pq}{(1-q)^3}=\frac{2q}{p^2}.
\]
So $\Exp[X^2]=\frac{2q}{p^2}+\frac1p$ and
$\Var(X)=\frac{2q}{p^2}+\frac1p-\frac1{p^2}=\frac{2q+p-1}{p^2}=\frac q{p^2}$.
\end{proof}

\begin{theorem}[Memorylessness]\label{thm:memoryless}
If $X\sim\mathrm{Geom}(p)$ then for non-negative integers $m,n$,
\[
\Prob(X>m+n\given X>m)=\Prob(X>n).
\]
Conversely, if $X$ takes values in $\{1,2,\dots\}$, has $\Prob(X>n)>0$ for all $n$,
and satisfies this identity, then $X$ is geometric.
\end{theorem}
\begin{proof}
We have $\Prob(X>n)=q^n$. Then
$\Prob(X>m+n\given X>m)=q^{m+n}/q^m=q^n$.
Conversely, let $G(n)=\Prob(X>n)$. The condition says $G(m+n)=G(m)G(n)$, so
$G(n)=G(1)^n$ by induction. Setting $q=G(1)=\Prob(X>1)$ gives $\Prob(X>n)=q^n$
and so $\Prob(X=n)=q^{n-1}-q^n=q^{n-1}(1-q)$, the geometric PMF with
$p=1-q$.
\end{proof}

In words: the process ``forgets'' its past. If you have already failed $m$ times,
the number of further trials is distributed as if you were starting over.

\begin{nuance}{two conventions for the geometric distribution}
Some books define the geometric variable as the number of \emph{failures} before
the first success, taking values $0,1,2,\dots$ with mean $q/p$. Always check the
convention. In this book geometric counts trials, starting at $1$.
\end{nuance}

\begin{definition}[Negative binomial]
$X$ has the \emph{negative binomial}$(r,p)$ distribution if it is the trial number
on which the $r$th success occurs:
\[
\Prob(X=k)=\binom{k-1}{r-1}p^rq^{k-r},\qquad k=r,r+1,\dots
\]
\end{definition}

\emph{Derivation.} The $r$th success is on trial $k$ exactly when trial $k$ is a
success and exactly $r-1$ of the first $k-1$ trials are successes. The latter
has probability $\binom{k-1}{r-1}p^{r-1}q^{k-r}$, and multiplying by the
independent success probability $p$ on trial $k$ gives the formula. With $r=1$ we
recover the geometric distribution.

\begin{theorem}\label{thm:nbmean}
If $X\sim\mathrm{NB}(r,p)$ then $\Exp[X]=r/p$.
\end{theorem}
\begin{proof}
Use $k\binom{k-1}{r-1}=r\binom kr$. Then
\[
\Exp[X]=\sum_{k\ge r}k\binom{k-1}{r-1}p^rq^{k-r}
=\frac rp\sum_{k\ge r}\binom kr p^{r+1}q^{k-r}.
\]
The terms in the last sum are $\Prob(Y=k+1)$ for $Y\sim\mathrm{NB}(r+1,p)$, since
$\binom{(k+1)-1}{(r+1)-1}p^{r+1}q^{(k+1)-(r+1)}=\binom kr p^{r+1}q^{k-r}$. The
PMF of $Y$ sums to $1$ over $k+1=r+1,r+2,\dots$, so the sum equals $1$.
\end{proof}

(The mean $r/p$ is expected: waiting for $r$ successes is $r$ consecutive waits,
each with mean $1/p$.)

\section{The Poisson distribution}

\begin{definition}[Poisson]
$X$ has the \emph{Poisson}$(\lambda)$ distribution, $\lambda>0$, if
\[
\Prob(X=k)=e^{-\lambda}\frac{\lambda^k}{k!},\qquad k=0,1,2,\dots
\]
\end{definition}

Total mass: $\sum_{k\ge0}e^{-\lambda}\lambda^k/k!=e^{-\lambda}e^\lambda=1$ by the
exponential series.

\begin{theorem}[Poisson limit of the binomial]\label{thm:poissonlimit}
Let $X_n\sim\mathrm{Bin}(n,p_n)$ with $np_n\to\lambda\in(0,\infty)$. Then for
every fixed $k\ge0$,
\[
\Prob(X_n=k)\longrightarrow e^{-\lambda}\frac{\lambda^k}{k!}.
\]
\end{theorem}
\begin{proof}
Write $\lambda_n=np_n\to\lambda$. Then
\[
\Prob(X_n=k)=\frac{n(n-1)\cdots(n-k+1)}{k!}\,\frac{\lambda_n^k}{n^k}
\Bigl(1-\frac{\lambda_n}n\Bigr)^{n-k}.
\]
As $n\to\infty$ with $k$ fixed: $\frac{n(n-1)\cdots(n-k+1)}{n^k}=
\prod_{j=0}^{k-1}(1-j/n)\to1$; $\lambda_n^k\to\lambda^k$;
$(1-\lambda_n/n)^{-k}\to1$; and $(1-\lambda_n/n)^n\to e^{-\lambda}$. For the last
limit, take logarithms: $n\ln(1-\lambda_n/n)=-\lambda_n+O(\lambda_n^2/n)\to-\lambda$,
using $\ln(1-x)=-x+O(x^2)$ as $x\to0$. Multiplying the limits gives the claim.
\end{proof}

\begin{theorem}[Moments of the Poisson]\label{thm:poissmoments}
If $X\sim\mathrm{Poisson}(\lambda)$ then $\Exp[X]=\Var(X)=\lambda$.
\end{theorem}
\begin{proof}
$\Exp[X]=\sum_{k\ge1}k\,e^{-\lambda}\frac{\lambda^k}{k!}=\lambda\sum_{k\ge1}
e^{-\lambda}\frac{\lambda^{k-1}}{(k-1)!}=\lambda$. Likewise
$\Exp[X(X-1)]=\sum_{k\ge2}e^{-\lambda}\frac{\lambda^k}{(k-2)!}=\lambda^2$, so
$\Var(X)=\lambda^2+\lambda-\lambda^2=\lambda$.
\end{proof}

The binomial limit explains the meaning: Poisson models the number of rare events
among many opportunities, where each event is unlikely but the expected total
$\lambda$ is moderate (misprints on a page, calls arriving at an exchange in a
minute, radioactive decays). It is also consistent with the moments: the binomial
has mean $np\to\lambda$ and variance $npq\to\lambda$.

\begin{nuance}{how good is the Poisson approximation?}
For $X=\sum_{i=1}^nI_i$ a sum of \emph{independent} Bernoulli$(p_i)$ indicators
with $\lambda=\sum p_i$, the distance between the law of $X$ and the
Poisson$(\lambda)$ law (in total variation) is at most $\sum_ip_i^2$ (Le Cam's
inequality; we cite it without proof). The approximation is good when every
individual $p_i$ is small, even if the $p_i$ differ. It can also be good under
weak dependence (the Chen--Stein method), as in the birthday problem for matching
pairs.
\end{nuance}

\begin{example}[Misprints]
A book has $500$ pages and $400$ misprints, each independently on a uniformly
random page. The number on a given page is $\mathrm{Bin}(400,1/500)$, with mean
$0.8$. Poisson$(0.8)$ gives $\Prob(\text{no misprint on page 1})\approx
e^{-0.8}\approx0.449$; the exact value is $(499/500)^{400}\approx0.4489$.
\end{example}

\section{The hypergeometric distribution}

An urn contains $N$ balls, $K$ of them red. Draw $n$ balls \emph{without
replacement}. The number $X$ of red balls drawn has the \emph{hypergeometric}
distribution:
\[
\Prob(X=k)=\frac{\binom Kk\binom{N-K}{n-k}}{\binom Nn},
\]
for integers $k$ with $\max(0,n-N+K)\le k\le\min(n,K)$. Counting argument: all
$\binom Nn$ subsets of size $n$ are equally likely; those with exactly $k$ red
balls are formed by choosing the red ones and the non-red ones separately. The
probabilities sum to $1$ by Vandermonde's identity (Theorem~\ref{thm:vander}).

\begin{theorem}[Mean of the hypergeometric]
$\Exp[X]=n\dfrac KN$.
\end{theorem}
\begin{proof}
Number the draws $1,\dots,n$ and let $I_j$ indicate that draw $j$ is red. Then
$X=\sum_jI_j$. By symmetry each draw is equally likely to be any of the $N$ balls,
so $\Prob(I_j=1)=K/N$, even though the draws are dependent. Linearity of
expectation gives $\Exp[X]=nK/N$.
\end{proof}

The variance $n\frac KN\bigl(1-\frac KN\bigr)\frac{N-n}{N-1}$ is derived in
Example~\ref{ex:hypervar} of the next chapter. Notice that it is the binomial
variance (drawing with replacement, $p=K/N$) times the \emph{finite population
correction} $\frac{N-n}{N-1}$.

\begin{proposition}[Binomial limit of the hypergeometric]
Let $N\to\infty$ with $K/N\to p\in(0,1)$ and $n$ fixed. Then the hypergeometric
PMF converges to the Bin$(n,p)$ PMF.
\end{proposition}
\begin{proof}
Write
$\Prob(X=k)=\binom nk\dfrac{K(K-1)\cdots(K-k+1)\cdot(N-K)\cdots(N-K-n+k+1)}
{N(N-1)\cdots(N-n+1)}$ (counting ordered sequences with a specified pattern of
red and non-red positions, times the $\binom nk$ patterns). The numerator has $k$
factors each $\sim pN$ and $n-k$ factors each $\sim(1-p)N$; the denominator has $n$
factors each $\sim N$. Hence the ratio tends to $p^k(1-p)^{n-k}$.
\end{proof}

\section{Summary}

\begin{center}
\begin{tabular}{@{}lccc@{}}
\toprule
Distribution & PMF & Mean & Variance\\
\midrule
Bernoulli$(p)$ & $p^k q^{1-k}$, $k\in\{0,1\}$ & $p$ & $pq$\\[0.3em]
Binomial$(n,p)$ & $\binom nk p^kq^{n-k}$ & $np$ & $npq$\\[0.3em]
Geometric$(p)$ & $q^{k-1}p$, $k\ge1$ & $1/p$ & $q/p^2$\\[0.3em]
NegBin$(r,p)$ & $\binom{k-1}{r-1}p^rq^{k-r}$, $k\ge r$ & $r/p$ & $rq/p^2$\\[0.3em]
Poisson$(\lambda)$ & $e^{-\lambda}\lambda^k/k!$ & $\lambda$ & $\lambda$\\[0.3em]
Hypergeom$(N,K,n)$ & $\binom Kk\binom{N-K}{n-k}/\binom Nn$ & $nK/N$ &
$n\frac KN\frac{N-K}N\frac{N-n}{N-1}$\\
\bottomrule
\end{tabular}
\end{center}

The negative binomial variance $rq/p^2$ is verified in Chapter~\ref{ch:joint}
(sum of $r$ independent geometrics).

\begin{interview}{recognising the distribution from the story}
Most exam problems say which distribution to use if you identify the story:
\emph{fixed number of independent trials, count successes}: binomial;
\emph{trials until first success}: geometric; \emph{trials until $r$th success}:
negative binomial; \emph{sampling without replacement from a finite population,
count a type}: hypergeometric; \emph{rare events in a large number of
opportunities}: Poisson. Always state which assumption (independence,
constant $p$) justifies the choice.
\end{interview}

\section*{Exercises}
\addcontentsline{toc}{section}{Exercises}

\begin{enumerate}
\item A fair coin is tossed $10$ times. Find the probabilities of exactly $5$
heads, at least $8$ heads, and the most likely number of heads.
\item A basketball player makes each free throw with probability $0.7$,
independently. What is the probability the first miss occurs on the fourth shot?
The expected number of shots until the first miss?
\item Prove: if $X\sim\mathrm{Poisson}(\lambda)$, then $\Prob(X=k+1)/\Prob(X=k)=
\lambda/(k+1)$, and deduce the mode is $\lfloor\lambda\rfloor$.
\item A call centre receives on average $3$ calls per minute (Poisson). Find the
probability of no calls in a given minute, and of more than $5$ calls.
\item Banach's matchbox: a person has two boxes with $n$ matches each, picking a
box at random for each use. When a box is found empty, the other contains $k$
matches. Show the probability is
$\binom{2n-k}{n}2^{-(2n-k)}$. (Count the draws carefully; there are two
symmetric cases.)
\item Show $\Exp[1/(1+X)]=\frac{1-q^{n+1}}{(n+1)p}$ for $X\sim\mathrm{Bin}(n,p)$.
\item In a lottery $6$ numbers are drawn from $49$. Find the probability that a
ticket with $6$ chosen numbers matches exactly $3$.
\item Prove the memoryless property for the geometric distribution in the form
$\Exp[X-m\given X>m]=\Exp[X]$.
\item (Challenge) Prove the Poisson limit theorem (Theorem~\ref{thm:poissonlimit})
for hypergeometric variables: if $N\to\infty$, $K\to\infty$, $n\to\infty$ with
$nK/N\to\lambda$ and $n/N\to0$, then the hypergeometric PMF converges to Poisson.
\end{enumerate}
